\section{The Two Cost Laws}
\label{sec:cost_laws}

With quality fixed at the exchange rate of \S3, we derive and verify the
per-event latency laws for the two architectures. The laws have
fundamentally different shapes: AR cost is a line in one variable
(output length); DG cost is a staircase in output length with a second,
steeper dependence on input length. Both laws are measured on the
pinned stack (vLLM~\vllmVersion, torch~\torchVersion, \gpuModel) with
run-to-run coefficient of variation below $0.3\%$, and both generalise
from BGL to HDFS within $1\%$.

\subsection{The AR cost law: linear in one variable}
\label{sec:cost_law_ar}

Autoregressive decoding generates one token at a time, each requiring a
full forward pass over the KV-cache. Cost is therefore linear in output
length and independent of how many canvases the output spans:
\begin{equation}
\label{eq:ar_law}
T_{\mathrm{AR}}(L_{\mathrm{out}}) \;=\; \arBase \;+\; \arSlope \,\cdot\, L_{\mathrm{out}}
\quad\text{(ms)}
\end{equation}
The intercept $\arBase$\,ms absorbs prompt prefill (amortised by prefix
caching across requests) and dispatch overhead; the slope
$\arSlope$\,ms/token is the per-token decode cost. We verify
Eq.~\ref{eq:ar_law} at two output lengths with five (resp.\ three)
repeats of $n{=}300$ events each on BGL.

\begin{table}[t]
\centering
\caption{AR-3B anchor measurements on the pinned stack. The linear law
$T = \arBase + \arSlope \cdot L$ predicts both points within $0.1\%$.
Run-to-run CV is below $0.06\%$.}
\label{tab:ar_anchor}
\begin{tabular}{lrrrr}
\toprule
$L_{\mathrm{out}}$ & Measured p50 (ms) & Predicted (ms) & Error & CV \\
\midrule
256 & 1616.0 & 1617.5 & $+0.09\%$ & $0.05\%$ \\
512 & 3225.1 & 3225.2 & $+0.00\%$ & $0.01\%$ \\
\bottomrule
\end{tabular}
\end{table}

Input length enters weakly. Sweeping the prompt from 992 to 6{,}172
tokens (by repeating the schema block) yields a slope of
$\arInputSlope$\,ms/prompt-token---less than $1\%$ of the per-token
decode slope. AR pays for the prompt once via prefix caching; each
generated token attends to the cached KV without recomputation.

The law generalises to HDFS: at $L{=}256$ the HDFS p50 is
1616.5\,ms---a $+0.03\%$ shift from BGL's 1616.0\,ms, well inside the
run-to-run CV.

\subsection{The DG cost law: staircase in two variables}
\label{sec:cost_law_dg}

Block-diffusion decoding generates tokens in fixed-size canvases of
$\canvasLength$ tokens. Within a canvas, all positions are denoised in
parallel at no additional sequential cost; crossing a canvas boundary
triggers a fresh denoising pass. Cost is therefore a staircase in
output length, not a line:
\begin{equation}
\label{eq:dg_law}
T_{\mathrm{DG}}(L_{\mathrm{in}},\, L_{\mathrm{out}}) \;=\;
\underbrace{\dgBase}_{\text{base (K=8)}} \;+\;
\underbrace{\dgInputSlope \cdot \Delta L_{\mathrm{in}}}_{\text{prefill}} \;+\;
\underbrace{\dgStep \cdot \Bigl(\!\Bigl\lceil \tfrac{L_{\mathrm{out}}}{\canvasLength}\Bigr\rceil - 1\Bigr)}_{\text{canvas staircase}}
\quad\text{(ms)}
\end{equation}
where $\Delta L_{\mathrm{in}} = L_{\mathrm{in}} - L_{\mathrm{in,ref}}$
is the deviation from the reference prompt ($\approx$1047 tokens). At
the reference prompt the prefill term vanishes and the law reduces to a
one-variable staircase with base $\dgBase$\,ms and step
$\dgStep$\,ms per additional canvas.

The staircase has sharp edges at exactly $\canvasLength$-token
boundaries and is flat between them. Table~\ref{tab:dg_edges} probes
both sides of the first two boundaries.

\begin{table}[t]
\centering
\caption{Canvas-boundary edges (DG, K=8, BGL). Crossing a 256-token
boundary costs $\sim$120\,ms; staying within a canvas costs $<$1\,ms.}
\label{tab:dg_edges}
\begin{tabular}{lrrll}
\toprule
$L_{\mathrm{out}}$ & p50 (ms) & $\lceil L/256\rceil$ & $\Delta$ from prev. & Edge? \\
\midrule
240 & 378.1 & 1 & --- & --- \\
256 & 378.4 & 1 & $+0.2$\,ms & no \\
264 & 500.0 & 2 & $+121.7$\,ms & \textbf{yes} \\
\midrule
496 & 500.0 & 2 & --- & --- \\
512 & 500.9 & 2 & $+0.9$\,ms & no \\
528 & 619.6 & 3 & $+118.7$\,ms & \textbf{yes} \\
\bottomrule
\end{tabular}
\end{table}

The step cost $\dgStep$\,ms is the median of seven canvas transitions
measured from $L{=}256$ to $L{=}2048$ (range 113.1--121.8\,ms, see
Table~\ref{tab:dg_staircase}). The variation is run-to-run noise, not a
systematic trend: the first extra canvas is not cheaper than later
ones, refuting an earlier hypothesis based on a K=32 vs K=8
configuration confound (\S0, discrepancy~\#2).

\begin{table}[t]
\centering
\caption{Full DG staircase on BGL (K=8, $n{=}200$ per point). Each row
is a canvas boundary; the step column is the p50 increment from the
previous row. Steps are uniform within run-to-run noise.}
\label{tab:dg_staircase}
\begin{tabular}{rrrr}
\toprule
$L_{\mathrm{out}}$ & canvases & p50 (ms) & step (ms) \\
\midrule
\phantom{0}256 & 1 & \phantom{0}378.4 & --- \\
\phantom{0}512 & 2 & \phantom{0}500.9 & 122.5 \\
\phantom{0}768 & 3 & \phantom{0}620.4 & 119.5 \\
1024 & 4 & \phantom{0}738.7 & 118.3 \\
1280 & 5 & \phantom{0}860.5 & 121.8 \\
1536 & 6 & \phantom{0}973.6 & 113.1 \\
1792 & 7 & 1091.2 & 117.5 \\
2048 & 8 & 1212.0 & 120.9 \\
\midrule
\multicolumn{2}{r}{median step} & & \dgStep \\
\bottomrule
\end{tabular}
\end{table}

Input length enters DG cost more steeply than AR. The same prompt sweep
yields a slope of $\dgInputSlope$\,ms/prompt-token for DG versus
$\arInputSlope$ for AR---a ratio of
$\dgInputSlope/\arInputSlope \approx 2.6\times$. Each denoising step
re-attends the full prompt, so the prefill cost is paid $K$ times
rather than once. This is the second variable that makes the DG law
two-dimensional.

Both variables generalise to HDFS. The AR law reproduces at $+0.03\%$
(\S\ref{sec:cost_law_ar}). The DG staircase step reproduces at
120.5\,ms on HDFS versus a BGL median of 118.5\,ms---within $1.7\%$.
The staircase shape is corpus-independent; only the prompt content
changes.

\subsection{Base decomposition: where the DG cost goes}
\label{sec:base_decomp}

The base $\dgBase$\,ms is not monolithic. Sweeping the denoising step
count $K \in \{2,4,8,16,32,64\}$ decomposes it into a $K$-independent
floor (prefill + fixed overhead) and a $K$-dependent denoising term:
\begin{equation}
\label{eq:ksweep}
T(K) \;=\; 296.1 \;+\; 2.57 \,\cdot\, K \quad\text{(ms, } R^2 = 0.49\text{)}
\end{equation}
The low $R^2$ reflects saturation above $K{=}16$: extra denoising steps
beyond convergence add negligible latency. Table~\ref{tab:ksweep}
reports the measured values and 9-field F1 at each $K$.

\begin{table}[t]
\centering
\caption{K-sweep: DG latency and quality vs denoising steps (BGL,
$L{=}256$, $n{=}200$). F1 collapses below $K{=}4$ and is flat above
$K{=}8$. Latency saturates above $K{=}16$.}
\label{tab:ksweep}
\begin{tabular}{rrrl}
\toprule
$K$ & p50 (ms) & 9-field F1 & regime \\
\midrule
\phantom{0}2 & \phantom{0}201.7 & 0.485 & collapse \\
\phantom{0}4 & \phantom{0}292.1 & 0.790 & transition \\
\phantom{0}8 & \phantom{0}378.6 & 0.853 & \textbf{operating point} \\
16 & \phantom{0}382.5 & 0.853 & flat \\
32 & \phantom{0}421.8 & 0.862 & flat \\
64 & \phantom{0}423.9 & 0.858 & saturated \\
\bottomrule
\end{tabular}
\end{table}

At the operating point $K{=}8$, the decomposition is:
\begin{equation}
\label{eq:decomp}
\underbrace{296\;\text{ms}}_{\substack{\text{prefill + fixed}\\\dgPrefillPct\%}}
\;+\;
\underbrace{20.6\;\text{ms}}_{\substack{\text{denoising loop}\\\dgDenoisePct\%}}
\;=\; 316.6\;\text{ms (fit)}
\end{equation}
The denoising loop---the part that makes the model ``diffusion''---is
$\dgDenoisePct\%$ of the fitted cost. The remaining $\dgPrefillPct\%$
is prefill and fixed dispatch overhead, shared with any
prompt-conditioned model. This is why halving $K$ from 8 to 4 saves only
$\sim$87\,ms ($23\%$) while collapsing F1 from 0.85 to 0.79: the cost
savings come from the small slice, the quality loss from the whole.

\subsection{The crossover}
\label{sec:crossover}

The two laws cross where $T_{\mathrm{AR}} = T_{\mathrm{DG}}$. At the
reference prompt ($\Delta L_{\mathrm{in}} = 0$) and the first canvas
($L_{\mathrm{out}} \leq \canvasLength$), the staircase reduces to its
floor and the crossover is analytic:
\begin{equation}
\label{eq:lcross}
L^* \;=\; \frac{\dgBase - \arBase}{\arSlope}
\;=\; \frac{378.8 - 7.27}{6.29}
\;=\; 59.1 \;\;\text{tokens}
\end{equation}
Below $L^* \approx \crossoverL$ output tokens, AR is faster; above it,
DG is faster---until the next canvas edge, where DG jumps by
$\dgStep$\,ms and the gap widens again.

The throughput crossover is the arrival rate at which DG's serial
single-event throughput saturates. With one event in flight, DG's
maximum throughput is the reciprocal of its base:
\begin{equation}
\label{eq:lambdacross}
\lambda^* \;=\; \frac{1000}{\dgBase}
\;=\; \frac{1000}{378.8}
\;=\; 2.6 \;\;\text{events/s}
\end{equation}
Above $\lambda^* \approx \crossoverLambda$\,ev/s, a queue forms in
front of the DG server. AR, which supports continuous batching, does
not saturate until $\sim$80\,ev/s (\S7).

Both crossovers are reproduced by the decision procedure's self-test
suite (\texttt{models/decision/test\_procedure.py}, 11/11 passing).
The test computes $L^*$ and $\lambda^*$ from the fitted constants
alone---no hardcoded values---confirming that the crossovers are
consequences of the laws, not tunable parameters.

\paragraph{Takeaway.} AR cost is a line through the origin with slope
$\arSlope$\,ms/token; DG cost is a staircase with floor $\dgBase$\,ms,
step $\dgStep$\,ms per $\canvasLength$-token canvas, and a prefill slope
$\dgInputSlope$\,ms/prompt-token that is $2.6\times$ steeper than AR's.
The lines cross at $L^* \approx \crossoverL$ tokens and $\lambda^*
\approx \crossoverLambda$\,ev/s. The denoising loop that distinguishes
DG from a cached-prefill model is $\dgDenoisePct\%$ of its own cost.
\S5 turns these laws into a deployment decision.
